\section{Policy extensions}\label{si:policy}

\paragraph{Reserved biocapacity share.} With a common reservation share $\sigma_f=\sigma_c=\sigma<1$, human-available supply is $\sigma B$ and the shortfall is $S=[E-\sigma B]_+$. The reservation lowers the demand at which shortfall-driven conversion begins; conversion is inactive while $E\le\sigma B$. With distinct shares, the corresponding condition is $E\le\sigma_fY_f+\sigma_cY_c$, which need not equal $\sigma B$. These are conditional results for the specified parameterization and do not establish robustness to disturbances or parameter uncertainty.

\paragraph{Typed floors and conversion stop.} Reserve $A_r\ge A_r^{\min}$, capital viability $A_c\ge\Acmin$, arable ceiling $A_f\le A_f^{\max}$, and $K\ge K_{\min}$ are separate constraints. A conversion stop triggered directly by $A_c=\Acmin$ preserves the capital-land floor. An $A_f$ ceiling constrains fast-land extent but does not by itself guarantee that floor unless the stop rule is linked to the capital floor and remaining-area constraints. In Eq.~(5), the conversion rate tends to zero as $A_c\to\Acmin$. An operator based only on aggregate $B$ cannot guarantee the typed capital-land floor because $B$ does not identify $A_c$.

\paragraph{Institutional signature.} Related work distinguishes protective from mobilizing institutional feedback \citep{Abaee2026a}. In this model, surplus-gated restoration is inactive during a shortfall, whereas deficit-gated restoration acts during it and can reduce provision further (Section~\ref{si:gates}).

\paragraph{Aggregation and typed feasibility.} Related transition-assessment work defines a linear compensatory operator and tests perfect substitutability against less-substitutable CES alternatives \citep[Sections~2 and 4 of][]{Abaee2026b}. The ECOMOD total $B$ is also a linear sum of land contributions. In the companion's separate finite-plan setting, the compensatory protocol accepts states for which no single available plan meets every component floor (Theorem~2 in \citet{Abaee2026b}); no positive-elasticity member is uniformly safe over the specified acceptance-gap region (Theorem~4(iii) in \citet{Abaee2026b}). Under pathwise feasibility, a decline below a typed floor during the transition remains a violation even if the endpoint later recovers. These results distinguish an aggregate account criterion from constraints on individual components.

\section{The one-stock comparator}\label{si:onestock}

This one-stock formulation is a structural comparator, not a reduction of the two-land model. Its state $A$ has comparator-specific stock units and is not the two-land state $A_c$. Let
\begin{align*}
G(A)&=\rho A\left(1-\frac{A}{A_{\max}}\right),\qquad
B(t)=b(t)A(t)+b_GG(A(t)),\qquad E(t)=eP(t);\\
\dot A(t)&=G\bigl(A(t-\tau_g)\bigr)-\frac{\pos{E(t)-\sigma b(t)A(t)}}{b_G};\\
\dot P(t)&=rP(t)\left(1-\frac{P(t)}{K(t-\tau_p)}\right),\qquad
K(t)=\max\!\left(\frac{B(t)}{e},K_{\min}\right);\\
\dot D(t)&=\pos{E(t)-B(t)}-\eta D(t),\qquad
b(t)=\bigl(b_0+T_b(t)\bigr)e^{-\alpha D(t)};\\
T_b(t)&=\omega_{Tb}\frac{\Delta b}{1+\exp[-\kappa(t-t_w)]}.
\end{align*}
Here $\sigma$ is a dimensionless human-available share and $\omega_{Tb}\in\{0,1\}$ switches the logistic technology increment on or off. $A$, $A_{\max}$, and the imposed lower bound $A_{\mathrm{ext}}$ use comparator-specific stock units; $A_{\max}$ is the logistic growth scale, not an imposed upper bound, and $G(A)$ is stock-unit\,yr$^{-1}$.  The coefficients $b$, $b_0$, $T_b$, and $\Delta b$ are gha per stock unit, while $b_G$ is gha\,yr per stock unit, so both terms in $B$ are annual account values in gha. $E$ and $B$ use the GFN/NFA convention, with the annual basis implicit; $P$, $K$, and $K_{\min}$ are in cap, and $e$ is gha\,cap$^{-1}$. The leaky state $D$ is in gha\,yr, so $\dot D$ and $\eta D$ are in gha; $\alpha$ is gha$^{-1}$\,yr$^{-1}$. Rates $\rho,r,\eta,\kappa$ are in yr$^{-1}$, and delays $\tau_g,\tau_p$ and the wave midpoint $t_w$ are in yr. The numerical implementation imposes $A\ge A_{\mathrm{ext}}$, $P,D\ge0$, and $K\ge K_{\min}$ and uses the exact positive-part operator. Delayed histories are constant at their initial values; RK4 integration uses linear interpolation of delayed states. The figure steps are $0.1$\,yr for the full trajectory (checked at $0.05$ and $0.025$\,yr), $0.01$\,yr for the fixed-demand mask (with a $0.02$\,yr shortfall sweep), $0.2$\,yr for the delay pair (checked at $0.1$\,yr), and $0.5$\,yr for the finite-grid classification.

With fixed $b$, $T_b=0$, $D=0$, and $\sigma=1$, interior equilibria with $0<A<A_{\max}$ and $B(A)/e>K_{\min}$ satisfy $E=B(A)$ and $P=B(A)/e$, forming a one-parameter family. The corresponding curve is
\[
B(A)=(b+b_G\rho)A-\frac{b_G\rho}{A_{\max}}A^2,\qquad
B'(A)=b+b_G\rho-2b_G\rho\frac{A}{A_{\max}}.
\]
An interior maximum on $0<A<A_{\max}$ requires $b_G\rho>b$. In the debt-free constant-parameter set used for the delay and basin illustrations, $b=0.5$, $b_G=0.8$, and $\rho=0.05$, so $b_G\rho=0.04<b$ and $B(A)$ is increasing on $0<A<A_{\max}$. In the full-trajectory illustration, by contrast, $b(t)$ varies with the technology wave and debt; no monotonicity or equilibrium-fold claim is made for that time-varying case.

At the debt-free reference equilibrium $A^*=0.8$ for the constant-parameter set $(b,b_G,\rho,A_{\max},e,r)=(0.5,0.8,0.05,1.2,0.55,0.02)$, $B^*=0.4107$\,gha and $P^*=B^*/e=0.7467$\,cap. The characteristic equation has a neutral root tangent to the equilibrium family and a positive real root. The latter is $0.5883$\,yr$^{-1}$ at $(\tau_g,\tau_p)=(0,0)$ and $0.6250$\,yr$^{-1}$ at $(30,25)$\,yr. These are comparator-specific local results, not a Hopf classification or a stability claim for the two-land system.

Figures~\ref{fig:si_onestock_full}--\ref{fig:si_onestock_basin} show conditional simulations of this comparator. In Fig.~\ref{fig:si_onestock_delay}, the same initial state approaches $A_{\max}$ for $\tau_g=10$\,yr and reaches the imposed floor for $\tau_g=30$\,yr when $\tau_p=25$\,yr; this pair illustrates dependence on the regeneration delay for one initial condition, not a universal critical delay. Figure~\ref{fig:si_onestock_basin} classifies a finite grid of initial conditions with $\tau_p$ held fixed. The proportions of grid cells are not probabilities over initial conditions or measures of a continuous basin.

\begin{figure}[htbp]
\centering
\includegraphics[width=0.96\textwidth]{si_onestock_full_trajectory}
\caption{Full delayed one-stock comparator with a logistic technology increment. Parameters: $A_0=1.0$ stock unit, $P_0=0.45$\,cap, $b_0=0.5$, $b_G=0.6$, $\rho=0.08$\,yr$^{-1}$, $A_{\max}=1.2$ stock units, $e=0.55$\,gha\,cap$^{-1}$, $r=0.02$\,yr$^{-1}$, $\eta=0.03$\,yr$^{-1}$, $\alpha=0.2$\,gha$^{-1}$\,yr$^{-1}$, $\tau_g=15$\,yr, $\tau_p=25$\,yr, $A_{\mathrm{ext}}=0.02$ stock units, $K_{\min}=10^{-4}$\,cap, $\Delta b=0.5$\,gha per stock unit, $\kappa=0.04$\,yr$^{-1}$, and $t_w=120$\,yr; $\sigma=1$, $\omega_{Tb}=1$, and $D_0=0$. The positive-part operator is exact. $B$ and $E$ are in gha, $A$ in comparator-specific stock units, $P$ in cap, and $D$ in gha\,yr. The trajectory reaches the imposed lower bound near $721$\,yr and subsequently recovers as modeled demand falls; the bound is not absorbing. This is a conditional model illustration, not an observation or a two-land result.}
\label{fig:si_onestock_full}
\end{figure}

\begin{figure}[htbp]
\centering
\includegraphics[width=0.96\textwidth]{si_onestock_mask_window}
\caption{One-stock masking example with fixed demand and no delays. Parameters: $A_0=1.0$ stock unit, $E=0.60$\,gha, $b_0=0.5$, $b_G=0.8$, $\rho=0.05$\,yr$^{-1}$, $A_{\max}=1.2$ stock units, $\eta=0.05$\,yr$^{-1}$, $\alpha=0.03$\,gha$^{-1}$\,yr$^{-1}$, $\Delta b=1.5$\,gha per stock unit, $\kappa=0.2$\,yr$^{-1}$, and $t_w=15$\,yr; $\sigma=1$, $\omega_{Tb}=1$, and $D_0=0$. Under the exact positive-part operator, the interval with $\dot B>0$ and $\dot A<0$ lasts about $3.0$\,yr; over it, $A$ falls from about $0.937$ to $0.867$ stock units while $B$ rises from about $0.575$ to $0.600$\,gha. Panel (b) shows the interval width across initial shortfalls $[E(0)-b(0)A_0]_+$ for this parameterization. These are conditional comparator outputs.}
\label{fig:si_onestock_mask}
\end{figure}

\begin{figure}[htbp]
\centering
\includegraphics[width=0.96\textwidth]{si_onestock_delay_examples}
\caption{Controlled regeneration-delay examples in the debt-free one-stock subsystem. Both runs use $A_0=1.0$ stock unit, $P_0=0.1$\,cap, $b=0.5$, $b_G=0.8$, $\rho=0.05$\,yr$^{-1}$, $A_{\max}=1.2$ stock units, $e=0.55$\,gha\,cap$^{-1}$, $r=0.02$\,yr$^{-1}$, $\tau_p=25$\,yr, $A_{\mathrm{ext}}=0.02$ stock units, $K_{\min}=10^{-6}$\,cap, and $\sigma=1$; only $\tau_g$ differs ($10$ versus $30$\,yr). $B$ is in gha, $P$ in cap, and $A$ in comparator-specific stock units. The selected trajectories do not determine a general critical delay.}
\label{fig:si_onestock_delay}
\end{figure}

\begin{figure}[htbp]
\centering
\includegraphics[width=0.96\textwidth]{si_onestock_basin}
\caption{Finite-grid initial-condition classification in the debt-free one-stock subsystem. The debt-free runs use the same parameter set as Fig.~\ref{fig:si_onestock_delay}, with $K_{\min}=10^{-6}$\,cap and $\tau_p=25$\,yr; $\tau_g=0$\,yr in the left panel and $30$\,yr in the right; the imposed lower bound is $A_{\mathrm{ext}}=0.02$ stock units. The grid is $A_0=0.10,0.20,\ldots,1.30$ stock units and $P_0=0.05,0.15,\ldots,1.55$\,cap; trajectories are integrated to $1200$\,yr with $\Delta t=0.5$\,yr. Green denotes a high endpoint $A(1200)\ge0.95A_{\max}$, red a low endpoint $A(1200)\le0.05$ stock units, and gray an intermediate endpoint. These display-grid classifications are not probabilities or continuous-basin measures.}
\label{fig:si_onestock_basin}
\end{figure}

\section{Controlled technology-wave experiment}\label{si:tbwave}

The matched experiment compares endogenous-demand runs with $T_b$ active or inactive, using $P_0=1.0$ or $1.2$\,cap. Both runs in each pair share $A_{tot}=2.5$\,ha, $A_{f0}=0.6$\,ha, $A_{c0}=0.9$\,ha, $q_0=0.9$, and $D_0=0.0$\,gha\,yr. Parameter values and numerical methods match; only $\omega_{Tb}$ differs. The common pre-run histories are constant at $q_0=0.9$ and the no-wave initial account value $B=0.560184$\,gha in both runs. Shared parameters are $b_{f0}=0.85$\,gha\,ha$^{-1}$, $b_c=0.05$\,gha\,ha$^{-1}$, $b_{G,c}=0.8$\,gha\,yr\,ha$^{-1}$, $\rho_c=0.08$\,yr$^{-1}$, $\alpha=0.03$\,gha$^{-1}$\,yr$^{-1}$, $\eta=0.05$\,yr$^{-1}$, $\eta_f=0.05$\,yr$^{-1}$, $\mu=0.06$\,yr$^{-1}$, $e=0.55$\,gha\,cap$^{-1}$, $r=0.02$\,yr$^{-1}$, $\tau_g=20$\,yr, $\tau_p=25$\,yr, $\tau_{conv}=1$\,yr, and $\chi_r=0.10$\,gha$^{-1}$\,yr$^{-1}$. The active-wave parameters are $\Delta b=1.0$\,gha\,ha$^{-1}$, $\kappa_w=0.1$\,yr$^{-1}$, and $t_{wave}=40$\,yr; the integration horizon is $250$\,yr. The logistic increment is nonzero at $t=0$ ($T_b(0)=0.017986$\,gha\,ha$^{-1}$), giving $B(0)=0.560184$\,gha without the wave and $0.570976$\,gha with it. The difference at $t=0$ arises from the logistic tail; the land, quality, population, and debt states are otherwise matched, and the pre-run $B$ history remains common. The wave reaches half its increment at $t_{wave}=40$\,yr.

Figure~\ref{fig:si_wave_debt} shows the $P_0=1.0$ pair. Its all-time debt maximum occurs after the wave midpoint in both runs: $D_{max}=0.008126$\,gha\,yr at $86.15$\,yr without the wave, versus $0.020205$\,gha\,yr at $115.60$\,yr with it. At $P_0=1.2$, by contrast, the all-time maximum occurs early and is lower with the wave, while the maximum over $t\ge t_{wave}$ is higher. Table~\ref{tab:si_wave_debt} reports both all-time maxima and maxima over $t\ge t_{wave}$. Because $t_{wave}$ is the logistic midpoint, the post-midpoint window includes the continued rise of the wave.

\begin{table}[htbp]
\centering\small
\caption{Matched technology-wave/no-wave debt maxima. Entries are $D_{max}$ in gha\,yr, with the time of the maximum in years in parentheses; the ratio is wave/no-wave. ``All'' is the maximum over $0\le t\le250$\,yr; $t\ge40$ is the post-midpoint maximum. The $P_0=1.0$ all-time peaks both occur after the midpoint.}
\label{tab:si_wave_debt}
\begin{tabularx}{\textwidth}{r l r l >{\centering\arraybackslash}X >{\centering\arraybackslash}X r}
\toprule
$P_0$ (cap) & quality update & $\Delta t$ (yr) & window & no wave $D_{max}$ (time) & wave $D_{max}$ (time) & ratio \\
\midrule
$1.0$ & delayed & $0.050$ & all & $0.008126\ (86.15)$ & $0.020205\ (115.60)$ & $2.487$ \\
$1.0$ & delayed & $0.025$ & all & $0.008105\ (86.20)$ & $0.020195\ (115.65)$ & $2.492$ \\
$1.0$ & current & $0.050$ & all & $0.008326\ (85.70)$ & $0.020168\ (115.60)$ & $2.422$ \\
$1.2$ & delayed & $0.050$ & all & $0.061218\ (1.65)$ & $0.051037\ (1.45)$ & $0.834$ \\
$1.2$ & delayed & $0.050$ & $t\ge40$ & $0.010854\ (81.70)$ & $0.025184\ (112.55)$ & $2.320$ \\
$1.2$ & delayed & $0.025$ & all & $0.060732\ (1.625)$ & $0.050582\ (1.475)$ & $0.833$ \\
$1.2$ & delayed & $0.025$ & $t\ge40$ & $0.010808\ (81.775)$ & $0.025165\ (112.625)$ & $2.328$ \\
\bottomrule
\end{tabularx}
\end{table}

For $P_0=1.0$, the all-time wave/no-wave debt-peak ratio is $2.487$ at $\Delta t=0.05$\,yr, $2.492$ at $0.025$\,yr, and $2.422$ under the current-quality specification at $0.05$\,yr. For $P_0=1.2$, the corresponding all-time ratios are $0.834$ and $0.833$, while the post-midpoint ratios are $2.320$ and $2.328$. Debt maxima change by less than $0.9\%$ under step halving. In the $P_0=1.0$ pair, $A_c$ does not fall below its initial value of $0.9$\,ha; the experiment therefore illustrates a conditional delayed debt rebound, not the capital-land drawdown in the frozen-demand baseline. The $P_0=1.2$ contrast shows that the all-time maximum and the post-midpoint maximum are distinct summaries. $D$ is the model's leaky time-integral of shortfall, not an observed NFA or physical-stock variable.

\begin{figure}[!htbp]
\centering
\includegraphics[width=0.92\textwidth]{fig_tb_wave_debt}
\caption{Matched endogenous-demand technology-wave/no-wave simulation at $P_0=1.0$\,cap. Parameters and initial conditions are listed in SI~\S S8; only $T_b$ is switched on or off. (a) The active logistic increment $T_b(t)$, with the dotted line at its midpoint $t_{wave}=40$\,yr; (b) annual account values $B$ and $E$ (gha); (c) population $P$ (cap); (d) leaky model debt $D$ (gha\,yr). Markers denote all-time maxima, which occur at $86.15$\,yr without the wave and $115.60$\,yr with it; peak debt is $0.008126$ versus $0.020205$\,gha\,yr. The plotted wave begins with its small logistic tail at $t=0$ and approaches $\Delta b=1.0$\,gha\,ha$^{-1}$; the dotted line marks the midpoint, not an end date.}
\label{fig:si_wave_debt}
\end{figure}

\section*{Supplementary references}
